% Options for packages loaded elsewhere
\PassOptionsToPackage{unicode}{hyperref}
\PassOptionsToPackage{hyphens}{url}
%
\documentclass[
  10pt,
  letterpaper,
]{article}
\usepackage{amsmath,amssymb}
\usepackage{iftex}
\ifPDFTeX
  \usepackage[T1]{fontenc}
  \usepackage[utf8]{inputenc}
  \usepackage{textcomp} % provide euro and other symbols
\else % if luatex or xetex
  \usepackage{unicode-math} % this also loads fontspec
  \defaultfontfeatures{Scale=MatchLowercase}
  \defaultfontfeatures[\rmfamily]{Ligatures=TeX,Scale=1}
\fi
\usepackage{lmodern}
\ifPDFTeX\else
  % xetex/luatex font selection
\fi
% Use upquote if available, for straight quotes in verbatim environments
\IfFileExists{upquote.sty}{\usepackage{upquote}}{}
\IfFileExists{microtype.sty}{% use microtype if available
  \usepackage[]{microtype}
  \UseMicrotypeSet[protrusion]{basicmath} % disable protrusion for tt fonts
}{}
\makeatletter
\@ifundefined{KOMAClassName}{% if non-KOMA class
  \IfFileExists{parskip.sty}{%
    \usepackage{parskip}
  }{% else
    \setlength{\parindent}{0pt}
    \setlength{\parskip}{6pt plus 2pt minus 1pt}}
}{% if KOMA class
  \KOMAoptions{parskip=half}}
\makeatother
\usepackage{xcolor}
\usepackage[margin=1in]{geometry}
\usepackage{color}
\usepackage{fancyvrb}
\newcommand{\VerbBar}{|}
\newcommand{\VERB}{\Verb[commandchars=\\\{\}]}
\DefineVerbatimEnvironment{Highlighting}{Verbatim}{commandchars=\\\{\}}
% Add ',fontsize=\small' for more characters per line
\newenvironment{Shaded}{}{}
\newcommand{\AlertTok}[1]{\textcolor[rgb]{1.00,0.00,0.00}{\textbf{#1}}}
\newcommand{\AnnotationTok}[1]{\textcolor[rgb]{0.38,0.63,0.69}{\textbf{\textit{#1}}}}
\newcommand{\AttributeTok}[1]{\textcolor[rgb]{0.49,0.56,0.16}{#1}}
\newcommand{\BaseNTok}[1]{\textcolor[rgb]{0.25,0.63,0.44}{#1}}
\newcommand{\BuiltInTok}[1]{\textcolor[rgb]{0.00,0.50,0.00}{#1}}
\newcommand{\CharTok}[1]{\textcolor[rgb]{0.25,0.44,0.63}{#1}}
\newcommand{\CommentTok}[1]{\textcolor[rgb]{0.38,0.63,0.69}{\textit{#1}}}
\newcommand{\CommentVarTok}[1]{\textcolor[rgb]{0.38,0.63,0.69}{\textbf{\textit{#1}}}}
\newcommand{\ConstantTok}[1]{\textcolor[rgb]{0.53,0.00,0.00}{#1}}
\newcommand{\ControlFlowTok}[1]{\textcolor[rgb]{0.00,0.44,0.13}{\textbf{#1}}}
\newcommand{\DataTypeTok}[1]{\textcolor[rgb]{0.56,0.13,0.00}{#1}}
\newcommand{\DecValTok}[1]{\textcolor[rgb]{0.25,0.63,0.44}{#1}}
\newcommand{\DocumentationTok}[1]{\textcolor[rgb]{0.73,0.13,0.13}{\textit{#1}}}
\newcommand{\ErrorTok}[1]{\textcolor[rgb]{1.00,0.00,0.00}{\textbf{#1}}}
\newcommand{\ExtensionTok}[1]{#1}
\newcommand{\FloatTok}[1]{\textcolor[rgb]{0.25,0.63,0.44}{#1}}
\newcommand{\FunctionTok}[1]{\textcolor[rgb]{0.02,0.16,0.49}{#1}}
\newcommand{\ImportTok}[1]{\textcolor[rgb]{0.00,0.50,0.00}{\textbf{#1}}}
\newcommand{\InformationTok}[1]{\textcolor[rgb]{0.38,0.63,0.69}{\textbf{\textit{#1}}}}
\newcommand{\KeywordTok}[1]{\textcolor[rgb]{0.00,0.44,0.13}{\textbf{#1}}}
\newcommand{\NormalTok}[1]{#1}
\newcommand{\OperatorTok}[1]{\textcolor[rgb]{0.40,0.40,0.40}{#1}}
\newcommand{\OtherTok}[1]{\textcolor[rgb]{0.00,0.44,0.13}{#1}}
\newcommand{\PreprocessorTok}[1]{\textcolor[rgb]{0.74,0.48,0.00}{#1}}
\newcommand{\RegionMarkerTok}[1]{#1}
\newcommand{\SpecialCharTok}[1]{\textcolor[rgb]{0.25,0.44,0.63}{#1}}
\newcommand{\SpecialStringTok}[1]{\textcolor[rgb]{0.73,0.40,0.53}{#1}}
\newcommand{\StringTok}[1]{\textcolor[rgb]{0.25,0.44,0.63}{#1}}
\newcommand{\VariableTok}[1]{\textcolor[rgb]{0.10,0.09,0.49}{#1}}
\newcommand{\VerbatimStringTok}[1]{\textcolor[rgb]{0.25,0.44,0.63}{#1}}
\newcommand{\WarningTok}[1]{\textcolor[rgb]{0.38,0.63,0.69}{\textbf{\textit{#1}}}}
\setlength{\emergencystretch}{3em} % prevent overfull lines
\providecommand{\tightlist}{%
  \setlength{\itemsep}{0pt}\setlength{\parskip}{0pt}}
\setcounter{secnumdepth}{-\maxdimen} % remove section numbering
\usepackage{xurl}
\setlength{\emergencystretch}{3em}
\usepackage{microtype}
\ifLuaTeX
  \usepackage{selnolig}  % disable illegal ligatures
\fi
\usepackage{bookmark}
\IfFileExists{xurl.sty}{\usepackage{xurl}}{} % add URL line breaks if available
\urlstyle{same}
\hypersetup{
  hidelinks,
  pdfcreator={LaTeX via pandoc}}

\author{}
\date{}

\begin{document}

\section{BA2-T1 - BAProposition shape and participation-role lower-bound
derivation}\label{ba2-t1---baproposition-shape-and-participation-role-lower-bound-derivation}

\textbf{Revision:} R1

\textbf{Status:} COMPLETED / PROVISIONAL CANDIDATE / BA2 NOT CLOSED

\textbf{Repository baseline reviewed:}
\texttt{3d8251328c77177375cccf1c51caa54b7473e21e}

\textbf{Phase:} BA2 - Relations and canonical action vocabulary

\textbf{BA0:} CLOSED

\textbf{BA1:} CLOSED

\textbf{BA3:} NOT STARTED

\subsection{1. Question and guardrail}\label{question-and-guardrail}

BA2-T1 asks only:

\begin{quote}
What is the minimum methodology-neutral structural shape required for
one \texttt{BAProposition} to express reusable project facts over
\texttt{BAReferent} identities without freezing an exhaustive predicate
vocabulary or importing a domain/method taxonomy?
\end{quote}

The test obeys the BA1 closure boundary:

\begin{Shaded}
\begin{Highlighting}[]
\NormalTok{BAReferent     accepted identity family}
\NormalTok{BAProposition  accepted identity family}
\end{Highlighting}
\end{Shaded}

A failure that requires a third semantic identity family would reopen
BA1. No such failure is found in this microstep.

BA2-T1 does not define provenance/source fields, views, ThreatForge
classes, STRIDE/DFD semantics, Findings or a complete relation/action
dictionary.

\subsection{2. Evidence pressure}\label{evidence-pressure}

\subsubsection{2.1 Documentation-layer
constraint}\label{documentation-layer-constraint}

The FunctionalRequirement metamodel makes normative prose semantically
primary and explicitly states that SPO references do not carry all
conditional, concurrency, lifecycle or failure semantics.

Its strong operational patterns include:

\begin{itemize}
\tightlist
\item
  condition -\textgreater{} required transition;
\item
  request -\textgreater{} processing responsibility -\textgreater{}
  observable result;
\item
  lifecycle event -\textgreater{} state change and failure behavior;
\item
  interaction -\textgreater{} output/correlation behavior;
\item
  concurrency condition -\textgreater{} atomicity/idempotency invariant.
\end{itemize}

Therefore BA2 cannot simply inherit documentation-layer SPO as the
complete analytical proposition model.

\subsubsection{2.2 Facial-access pressure}\label{facial-access-pressure}

\texttt{FR-3.4} states one delivery meaning with at least these
distinguishable semantic participants/conditions:

\begin{itemize}
\tightlist
\item
  \texttt{CameraSubsystem};
\item
  \texttt{RecognitionProcessor};
\item
  \texttt{RecognitionCapture};
\item
  \texttt{RecognitionRequest} correlation;
\item
  incomplete delivery must not be represented as success.
\end{itemize}

The same delivery meaning is later constrained by confidentiality,
integrity and authorized-origin requirements and survives some mutations
while being retired by the local-recognition mutation.

\subsubsection{2.3 Order-fulfillment
pressure}\label{order-fulfillment-pressure}

The independent order corpus adds propositions that combine:

\begin{itemize}
\tightlist
\item
  executor/capability;
\item
  request and result;
\item
  \texttt{OrderEvaluation} correlation;
\item
  \texttt{Reservation} lifecycle;
\item
  all-lines atomicity;
\item
  competing concurrency;
\item
  idempotency;
\item
  compensation;
\item
  authorized/declined/indeterminate outcomes;
\item
  physical handoff reuse across inventory and payment behavior.
\end{itemize}

The corpus explicitly says its actor/data/store/contract probe kinds are
authoring substrate only and that final Base Analysis types are
deferred.

\subsection{3. Alternative A - universal binary
SPO}\label{alternative-a---universal-binary-spo}

Candidate:

\begin{Shaded}
\begin{Highlighting}[]
\NormalTok{subject {-}{-} predicate {-}{-}\textgreater{} object}
\end{Highlighting}
\end{Shaded}

\subsubsection{Test}\label{test}

Binary edges can represent simple facts, and they remain useful
projections. The problem is using them as the \textbf{only} normative
proposition shape.

For \texttt{FR-3.4}, four binary edges could encode source, destination,
information and correlation only if there is a common semantic anchor
that says they belong to the same delivery assertion. If the anchor is
the \texttt{BAProposition}, the binary-only claim has already expanded
into a role-bearing n-ary assertion. If the anchor is a synthetic
\texttt{BAReferent} created for every relation occurrence, the model
violates the BA1 distinction by using referent identity as a generic
proposition proxy.

The order corpus makes the problem stronger: atomicity, concurrency and
conditional result semantics apply to the joint configuration of
multiple participants. Splitting them into unrelated triples requires a
separate grouping/scope mechanism before the original claim can be
reconstructed.

\subsubsection{Disposition}\label{disposition}

\textbf{REJECTED AS UNIVERSAL SHAPE.} Binary relations remain valid
special cases or projections of a more general proposition structure.

\subsection{4. Alternative B - one proposition with anonymous ordered
participants}\label{alternative-b---one-proposition-with-anonymous-ordered-participants}

Candidate:

\begin{Shaded}
\begin{Highlighting}[]
\NormalTok{operator(participant1, participant2, participant3, ...)}
\end{Highlighting}
\end{Shaded}

\subsubsection{Test}\label{test-1}

Positional ordering is insufficient because the same referent can occupy
different semantic roles in different propositions, and method consumers
should not depend on undocumented argument positions.

For example, \texttt{RecognitionCapture} may be transferred information
in one assertion and the object of an acceptance/validity condition in
another. \texttt{PaymentProvider} may be a destination/provider
participant in one proposition and an external authority/context in
another.

\subsubsection{Disposition}\label{disposition-1}

\textbf{REJECTED.} Participation roles must be explicit.

\subsection{5. Alternative C - fixed domain-specific
slots}\label{alternative-c---fixed-domain-specific-slots}

Candidate examples:

\begin{Shaded}
\begin{Highlighting}[]
\NormalTok{source / process / destination / data / store / trustBoundary}
\end{Highlighting}
\end{Shaded}

or another fixed architecture/threat-method slot set.

\subsubsection{Test}\label{test-2}

Facial access can be made to look DFD-like, but the order corpus needs
request/result, correlation, reservation lifecycle, responsibility,
atomicity, compensation and physical handoff semantics that do not fit
one fixed DFD tuple without distortion.

A domain-specific slot set would also violate the closed method-neutral
boundary.

\subsubsection{Disposition}\label{disposition-2}

\textbf{REJECTED AS UNIVERSAL CORE.} Domain/method projections may
derive such slots later.

\subsection{6. Lower bound 1 - explicit semantic
operator}\label{lower-bound-1---explicit-semantic-operator}

A set of participants and roles is still ambiguous unless the
proposition says what semantic assertion connects them.

Therefore every proposition needs an explicit \textbf{method-neutral
semantic operator key}.

This microstep does not accept a verb list. It accepts only the
structural responsibility:

\begin{Shaded}
\begin{Highlighting}[]
\NormalTok{BAProposition.semanticOperator = exactly 1 explicit semantic key}
\end{Highlighting}
\end{Shaded}

The key must eventually be stable enough for mechanical selection and
comparison, but whether canonical operators are modeled as controlled
symbols, registry entries or another representation remains later BA2
work.

No \texttt{Relation} or \texttt{Action} BAE identity family is created.

\subsection{7. Lower bound 2 - role-bound n-ary
participation}\label{lower-bound-2---role-bound-n-ary-participation}

A BAProposition must support one or more role-bound semantic terms:

\begin{Shaded}
\begin{Highlighting}[]
\NormalTok{participation 1..*}
\NormalTok{    role 1}
\NormalTok{    term 1}
\end{Highlighting}
\end{Shaded}

At least one term must be a \texttt{BAReferent}, preserving the BA1
definition that a proposition is an assertion about referent identities.

A term may be a typed local value when cross-proposition semantic
identity is not required. If a condition/state/value must itself be
reused, compared or constrained across propositions, BA1 already
provides the promotion path: represent that project meaning as a
\texttt{BAReferent}.

This form permits unary, binary and higher-arity propositions. N-ary
support is therefore a \textbf{capability}, not a requirement that every
proposition have more than two participants.

\subsubsection{Disposition}\label{disposition-3}

\textbf{FORCED CANDIDATE.}

\subsection{8. Lower bound 3 - scoped semantic
modifiers}\label{lower-bound-3---scoped-semantic-modifiers}

The reviewed clauses show that the truth/meaning of an assertion may
depend on:

\begin{itemize}
\tightlist
\item
  condition/applicability;
\item
  state;
\item
  failure/completion rule;
\item
  realization;
\item
  atomicity/concurrency;
\item
  idempotency;
\item
  polarity/negative outcome.
\end{itemize}

Treating this as an untyped prose suffix would defeat the purpose of
Base Analysis. But forcing every local qualifier to become a
\texttt{BAReferent} would over-reify project meaning.

BA2-T1 therefore accepts a weaker structural requirement:

\begin{quote}
a BAProposition must support explicit, proposition-scoped semantic
modifiers without requiring those modifiers to become BAE identities.
\end{quote}

The exact modifier schema is intentionally \textbf{OPEN}. BA2-T2 must
test whether modifiers are best represented through role-bound
values/referents, normalized qualifier records, separate propositions
about reusable referents, or a constrained combination.

\subsubsection{Disposition}\label{disposition-4}

\textbf{REQUIRED CAPABILITY / MATERIAL FORM OPEN.}

\subsection{9. Referent classification versus participation
role}\label{referent-classification-versus-participation-role}

The BA1 closure allowed method-neutral semantic classifications without
creating subtype roots. BA2-T1 asks whether those classifications can
instead be reconstructed from proposition roles.

They cannot reliably be collapsed.

Role is contextual; classification is reusable semantic characterization
of the referent. \texttt{RecognitionCapture} does not cease to be
information/resource-like project meaning because one proposition uses
it as an acceptance object. \texttt{PaymentProviderContract} does not
become a destination type merely because an interaction proposition
points toward it/provider context.

If a method-neutral consumer must know reusable semantic kind but Base
Analysis omits it, the consumer would have to reread raw prose or invent
its own taxonomy, weakening shared semantics.

\subsubsection{Disposition}\label{disposition-5}

\begin{Shaded}
\begin{Highlighting}[]
\NormalTok{classification derived only from roles         REJECTED}
\NormalTok{explicit reusable classification capability     SUPPORTED}
\NormalTok{exact classification vocabulary/cardinality     DEFERRED}
\end{Highlighting}
\end{Shaded}

This does not reopen BA1: classification remains cheaper than a
first-class split.

\subsection{10. Behavior/action identity
test}\label{behavioraction-identity-test}

BA1 already decided that a behavior/event may be a \texttt{BAReferent}
when independently reused. BA2-T1 tests whether proposition structure
now forces a dedicated \texttt{Action} family.

It does not.

A delivery behavior that must be referenced by multiple constraints may
be a BAReferent. A proposition's semantic operator still expresses the
assertion currently being made. These are different responsibilities:

\begin{Shaded}
\begin{Highlighting}[]
\NormalTok{reusable project behavior meaning  {-}\textgreater{} BAReferent}
\NormalTok{asserted semantic relation/action  {-}\textgreater{} BAProposition operator + participation}
\end{Highlighting}
\end{Shaded}

\subsubsection{Disposition}\label{disposition-6}

\textbf{NO BA1 REOPEN.} No third identity family is forced.

\subsection{11. Concrete facial-access
rendering}\label{concrete-facial-access-rendering}

The following labels are illustrative only:

\begin{Shaded}
\begin{Highlighting}[]
\NormalTok{P{-}FR34{-}DELIVERY}
\NormalTok{semanticOperator: transfer}
\NormalTok{participation:}
\NormalTok{  source              {-}\textgreater{} CameraSubsystem}
\NormalTok{  destination         {-}\textgreater{} RecognitionProcessor}
\NormalTok{  information         {-}\textgreater{} RecognitionCapture}
\NormalTok{  correlationContext  {-}\textgreater{} RecognitionRequest}
\NormalTok{scopedModifier:}
\NormalTok{  completion          {-}\textgreater{} incomplete delivery != successful completion}
\end{Highlighting}
\end{Shaded}

This is already richer than SPO but still smaller than a domain-specific
DFD ontology.

If \texttt{FR-3.4\ delivery} is independently reused by specialized
constraints, its project meaning may additionally be represented as a
BAReferent and participate explicitly in those propositions. The
assertion identity \texttt{P-FR34-DELIVERY} does not become the
project-semantic target.

\subsection{12. Concrete order-fulfillment
rendering}\label{concrete-order-fulfillment-rendering}

Consider the atomic reservation requirement. One proposition must be
able to expose, without synthetic document/tool types:

\begin{Shaded}
\begin{Highlighting}[]
\NormalTok{semanticOperator: reserve / allocate   [illustrative]}
\NormalTok{participation:}
\NormalTok{  executor     {-}\textgreater{} InventoryService}
\NormalTok{  request      {-}\textgreater{} ReservationRequest}
\NormalTok{  correlation  {-}\textgreater{} OrderEvaluation}
\NormalTok{  created      {-}\textgreater{} Reservation}
\NormalTok{  result       {-}\textgreater{} ReservationResult}
\NormalTok{  affected     {-}\textgreater{} InventoryBalance}
\NormalTok{scopedModifier:}
\NormalTok{  atomicity    {-}\textgreater{} all requested lines or none}
\NormalTok{  concurrency  {-}\textgreater{} competing reservation cannot consume frozen quantity}
\end{Highlighting}
\end{Shaded}

The exact canonical roles and operator names are not accepted here. The
example demonstrates only why the structure must be role-bound and able
to carry scoped semantics.

\subsection{13. Bounded consumer check}\label{bounded-consumer-check}

The prior bounded STRIDE-oriented consumer needs behavior, source,
destination, information, correlation, responsibility/externality,
realization, failure and constraints.

BA2-T1's candidate can expose those semantics without importing STRIDE
categories:

\begin{itemize}
\tightlist
\item
  referents carry stable project-semantic identity and reusable
  classification where needed;
\item
  propositions carry explicit operator semantics;
\item
  role bindings expose participation mechanically;
\item
  scoped modifiers preserve conditions/failure semantics;
\item
  method-specific categories remain in the consumer.
\end{itemize}

\subsubsection{Result}\label{result}

\textbf{STRUCTURAL PASS / VOCABULARY NOT YET CLOSED.}

\subsection{14. Falsification summary}\label{falsification-summary}

\begin{itemize}
\tightlist
\item
  Pure binary SPO as universal core: \textbf{REJECTED}.
\item
  Binary proposition as a valid special case: \textbf{SUPPORTED}.
\item
  Anonymous positional n-ary tuple: \textbf{REJECTED}.
\item
  Explicit role-bound n-ary participation: \textbf{FORCED CANDIDATE}.
\item
  Fixed domain/method slot model: \textbf{REJECTED}.
\item
  Explicit semantic operator key: \textbf{FORCED CANDIDATE}.
\item
  Exhaustive operator list: \textbf{DEFERRED}.
\item
  Scoped modifier capability: \textbf{REQUIRED / ENCODING OPEN}.
\item
  Free-text modifier bag: \textbf{REJECTED}.
\item
  Participation/Role/Qualifier as BAE identity families: \textbf{NOT
  FORCED}.
\item
  Referent classification reconstructed only from roles:
  \textbf{REJECTED}.
\item
  Explicit method-neutral classification capability: \textbf{SUPPORTED /
  VOCABULARY OPEN}.
\item
  Dedicated Action/Event/Relation identity family: \textbf{NOT FORCED}.
\item
  BA1 reopen: \textbf{NOT TRIGGERED}.
\end{itemize}

\subsection{15. Provisional candidate}\label{provisional-candidate}

\begin{Shaded}
\begin{Highlighting}[]
\NormalTok{BAProposition}
\NormalTok{|{-} identity                         [BA1 CLOSED]}
\NormalTok{|{-} semanticOperator                1}
\NormalTok{|{-} participation                   1..*}
\NormalTok{|    |{-} role                       1}
\NormalTok{|    \textasciigrave{}{-} term                       1}
\NormalTok{\textasciigrave{}{-} scopedModifier                  0..* [capability required; exact shape open]}
\end{Highlighting}
\end{Shaded}

Additional constraints:

\begin{enumerate}
\def\labelenumi{\arabic{enumi}.}
\tightlist
\item
  at least one participation term is a \texttt{BAReferent};
\item
  n-ary participation is supported but not mandatory for every
  proposition;
\item
  role semantics are explicit, not positional;
\item
  semantic operator and role vocabularies are method-neutral and still
  open;
\item
  reusable referent classification is separate from contextual
  participation role;
\item
  no new BAE identity family is introduced.
\end{enumerate}

\subsection{16. What BA2-T1 does not
decide}\label{what-ba2-t1-does-not-decide}

BA2-T1 does not decide:

\begin{itemize}
\tightlist
\item
  final operator vocabulary;
\item
  final participation-role vocabulary;
\item
  operator-role compatibility matrix;
\item
  role cardinalities;
\item
  exact qualifier/modifier representation;
\item
  logical composition/negation expression language;
\item
  canonical classification vocabulary;
\item
  provenance/source fields;
\item
  views/projections;
\item
  lexical synonyms;
\item
  threat-method mappings.
\end{itemize}

\subsection{17. Next authorized
microstep}\label{next-authorized-microstep}

Execute only:

\begin{quote}
\textbf{BA2-T2 - semantic operator, participation-role and
scoped-modifier vocabulary pressure test.}
\end{quote}

BA2-T2 must attempt to derive the smallest reusable method-neutral
semantic vocabulary from the two corpora and bounded consumer while
keeping lexical synonyms, method-specific categories and provenance out
of scope.

Do not close BA2 merely because the structural lower bound survives this
trial.

\end{document}
